\subsection{Which Genres Anchor Complex Tastes? Demographic Concentration}
\label{sec:genre-specificity}

Beyond the aggregate level, genre-specific disagreement reveals which demographic groups hold complex tastes concentrated in particular repertoires. The educated and affluent do not simply hold \emph{more} disagreement overall; they concentrate it in specific cultural objects---notably prestige genres like opera and classical music. Respondents whose parents hold college degrees show the strongest genre-specific disagreement in opera ($r = 0.15$ on Dimension 2, $r = 0.14$ on Dimension 1), followed by indie rock, classical, and Broadway/musicals. This pattern suggests that educational capital structures not just \emph{whether} one's aspects disagree, but \emph{where} that disagreement concentrates. Educated respondents navigate these prestige forms through contradictory positions: liking but not valuing, listening without granting high status, or valuing without regular engagement.

Subjective class position cuts a different path. Working-class and lower-class respondents concentrate their genre-specific disagreement in rap and hip-hop ($r = -0.13$), the inverse of the education pattern. These genres capture working-class guilty pleasures and distanced consumption: respondents like rap but withhold status attribution, or listen while maintaining an ironic distance from the genre's valuations. Higher-class respondents show more disagreement in Broadway/musicals and big band ($r = 0.13$ and $r = 0.11$), genres they are more likely to encounter and evaluate.

Age creates the sharpest genre-specific patterns. Younger respondents hold substantially more genre-specific disagreement in rap and hip-hop ($r = -0.43$), indie rock ($r = -0.34$), and dance/club music ($r = -0.24$)---the popular genres that define contemporary youth culture. Their disagreement reflects conflicted engagement: they may like these genres intensely but resist full identification with them, or listen widely while selectively evaluating. Older respondents, by contrast, concentrate disagreement in big band and jazz ($r = 0.19$ and modest positive effects), genres encountered less frequently and navigated with less fluency.

Gender differences emerge primarily in genres coded by masculine cultural expression. Men hold significantly more genre-specific disagreement in big band ($r = -0.12$) and heavy metal ($r = -0.10$)---both male-coded genres where complex configurations are more prevalent. Women show concentrated disagreement in hymns and gospel ($r = 0.09$), suggesting navigation between personal liking and collective religious or cultural valuations.

Race and ethnicity show perhaps the clearest genre-specific patterns. Non-white respondents, especially Hispanic or Latino, concentrate genre-specific disagreement heavily in Latin/salsa music ($r = 0.31$ on Dimension 1, $r = 0.18$ on Dimension 2)---their own cultural form, navigated through complex relationships between personal liking, listening patterns, and broader status valuations. Rap and hip-hop show similar racialized disagreement ($r = 0.21$), as do dance and club music ($r = 0.20$ and $r = 0.15$). These genres cluster where cultural participation and evaluative distance intersect most acutely: respondents of color holding complex, often contradictory positions toward genres that represent both personal cultural identity and contested social status.

Figure~\ref{fig:omega-demographics} maps these patterns visually across eight demographically interesting genres. The genre-based component thus becomes a window onto how demographic position shapes not just taste position but taste \emph{structure}---the specific locations where respondents experience their aspects pulling apart. Social location does not uniformly inflate disagreement; it localizes it, channeling complex tastes through particular cultural objects that carry demographic and social meaning.

\begin{figure}[htbp]
\centering
\includegraphics[width=.95\linewidth]{fig-omega-demographics}
\caption{Which demographic groups hold genre-specific disagreement? Lollipop chart shows correlations between genre-based disagreement ($\omega$) and five demographic variables (parental education, subjective class, female gender, age in years, race/ethnicity), across the three dimensions of complex taste space. Each point represents the correlation for one genre; closer to the center indicates weaker effect. Positive correlations mean the demographic variable is associated with more genre-specific disagreement in that genre; negative correlations mean the opposite.}
\label{fig:omega-demographics}
\end{figure}
